% !PRINTED 317
\weberpage{317}

Um dies einzusehen, bezeichnen wir in der früheren Bezeichnungsweise 1, 2, 3 mit $(\xi)$, $(\eta)$, $(\zeta)$; da diese drei Grössen ein Tripel bilden, so giebt es [nach \S.\,81, (5)] ein System $(\alpha)$, so dass
\weq{4}{\Sigma\alpha\xi \equiv 0, \quad \Sigma\alpha\eta \equiv 0, \quad \Sigma\alpha\zeta \equiv 0 \ (\mathrm{mod}\ 2)}
ist. Wenden wir nun eine Substitution aus $Q$ an:
\[
\begin{pmatrix} \xi_1, & \xi_2, & \xi_3 \\ \xi_1', & \xi_2', & \xi_3' \end{pmatrix},
\]
worin $\xi_i' = \xi_i + h_i$ sein mag, so folgt aus (4):
\[
\Sigma\alpha\xi' \equiv \Sigma\alpha\eta' \equiv \Sigma\alpha\zeta' \equiv \Sigma\alpha h,
\]
d.\,h. die $(\xi')$, $(\eta')$, $(\zeta')$ bilden das zu $(\alpha)$ gehörige Tripel, wenn $(h)$ selbst zu diesem Tripel gehört, oder im anderen Falle ist $(h)$, $(\xi')$, $(\eta')$, $(\zeta')$, das zu $(\alpha)$ gehörige Quadrupel.

Bezeichnet also nun wieder $k$ irgend eine Substitution aus $P$ und $\gamma$ eine Substitution aus $Q$, durch die das Quadrupel 0, 1, 2, 3 in 4, 5, 6, 7 übergeht, so können wir die Elemente $\tau'$, $\tau''$ aus $T_{24}$ und $\gamma'$, $\gamma''$ aus $Q$ so bestimmen, dass
\[
k = \tau'\gamma' = \gamma^{-1}\tau''\gamma'',
\]
wodurch, wie oben, bewiesen wird, dass die beiden Systeme (3) durch $k$ imprimitiv verbunden sind.

Nach Adjunction einer Quadratwurzel zum Rationalitätsbereiche sind dann die Grössen 0, 1, 2, 3 die Wurzeln einer biquadratischen Gleichung.

Durch diese Betrachtungen wird für primitive Gleichungen 8\textsuperscript{ten} Grades auch die Möglichkeit ausgeschlossen, dass eine der sieben Grössen $\Psi_{\xi_1,\xi_2,\xi_3}$ [\S.\,80, (1)] verschwinde. Diese Functionen sind nämlich rationale Functionen der Wurzeln $X$; wenn also eine von ihnen verschwindet, so kann auf die rationale Gleichung $\Psi = 0$ jede Permutation der Gruppe der Gleichung angewandt werden. Daraus ergiebt sich nach dem Theorem \S.\,80, 2., dass entweder alle sieben Functionen $\Psi$ verschwinden müssen, in welchem Falle die Wurzeln rational wären, oder dass die Substitutionen $\varrho^{-1}$, auf die $(\xi)$ angewandt, eine intransitive Gruppe bilden müssen. Der Grad der Gruppe könnte dann nicht durch 7 theilbar sein. Nun ist die Gruppe der $\varrho^{-1}$ isomorph mit der Gruppe $T$ der $\sigma$ (\S.\,37, 8.), und daher ist auch der Grad von $T$ nicht durch 7 theilbar, und folglich ist
%CONT>>
